\subsection{An Enriques component in a zero-Euler torus filling}

\begin{statement}{C.2.UN.25}{Proved}{An essential Enriques surface occurs in a compact zero-Euler torus filling}
There is a smooth compact complex threefold over
$\PP^1$ with connected two-torus general fibre and
a simple-normal-crossing special support of Euler
number zero containing an unchanged smooth Enriques
surface as an irreducible component.
Before six explicit vertical modifications, its full
central fibre is
\[
 2M+2A+4Q_1+4Q_2+E_1+E_2+E_3+E_4,                 \tag{C.EN.1}
\]
where $M$ is Enriques, $A\simeq\PP^1\times\PP^1$,
$Q_1,Q_2$ are primary Hopf surfaces, and every
$E_i$ is the Hirzebruch surface $\mathbb F_2$.
That reduced support has Euler number $12$.
All six modifications can be made disjoint from $M$.
The Enriques component is therefore essential in the
sense of C.2.UN.14.
The reduced supports before and after balancing have
respectively
\[
 \begin{array}{c|c|c}
 & (b_0,b_1,b_2,b_3,b_4)&(h^0(\OO),h^1(\OO),h^2(\OO))\\ \hline
 S&(1,1,8,4,8)&(1,1,0)\\
 S_{\rm bal}&(1,1,20,40,20)&(1,1,0).
 \end{array}                                             \tag{C.EN.18}
\]
\end{statement}
\begin{proof}
We first give a compact base with specified local
coordinates. Take the Weierstrass elliptic surface
over $\PP^1_T$
\[
 y^2=x^3+(T^4-3)x+(T^6+T+2).                       \tag{C.EN.2}
\]
The coefficients are sections of $\OO(4)$ and
$\OO(6)$. Its discriminant, up to the usual nonzero
constant factor, is
\[
 \begin{split}
 \Delta(T)={}&31T^{12}-36T^8+54T^7+108T^6\\
             &+108T^4+27T^2+108T .
 \end{split}                                       \tag{C.EN.3}
\]
The exact Euclidean algorithm gives
$\gcd(\Delta,\Delta')=1$; also
$\Delta'(0)=108$ and its leading coefficient is
$31$. These checks are saved in
\texttt{code/E\_enriques\_filling.py}.
Thus every singular fibre is nodal of type $I_1$,
the fibre at zero is nodal, and infinity is smooth.
The total Weierstrass surface is smooth.

Here is an explicit chart at the node of the
zero fibre. Put $\xi=x-1$ and
$U=1+T^3(1+\xi)+T^5$. On the surface,
\[
 X=\frac{y-\xi\sqrt{3+\xi}}{U},\qquad
 Y=y+\xi\sqrt{3+\xi}
 \quad\Longrightarrow\quad XY=T .                  \tag{C.EN.4}
\]
The coordinate Jacobian at $\xi=y=T=0$ is
nonzero, so these are actual local coordinates.
Make the quadratic base change $T=t^2$ and
resolve $XY=t^2$. One resolution chart is
\[
       (X,Y,t)=(u,uv^2,uv),\qquad
       h(u,v)=(u,-v),                              \tag{C.EN.5}
\]
where $h$ is the deck involution. The other chart
interchanges $X,Y$ and replaces $v$ by $v^{-1}$.
The central fibre of the smooth base surface $V_0$
is $R_0+R_1$ of type $I_2$. Its strict nodal
normalization $R_0\simeq\PP^1$ is fixed pointwise
by $h$. The exceptional curve $R_1\simeq\PP^1$
has coordinate $z=v$, with $h(z)=-z$, and
meets $R_0$ exactly at $z=0,\infty$.

Blow up the four points $z=1,2,-1,-2$ of $R_1$.
Write their exceptional curves as
$D_1,D_2,D_{-1},D_{-2}$, and denote the resulting
compact surface by $V$. The involution lifts
and exchanges $D_j,D_{-j}$. Set
\[
 D=D_1+D_2-D_{-1}-D_{-2},\qquad L=\OO_V(D).
                                                               \tag{C.EN.6}
\]
The canonical meromorphic section $\sigma$ of
$L$ supplies the section-preserving identification
$h^*L=L^{-1}$. Its square is the identity.

Apply the explicit signed root construction of
C.2.UN.18 with order $m=2$ at these four disjoint
curves. Thus on the complement of the negative
curves use $\sigma=a^2\ell$, with $\ell\in L^\times$;
on the complement of the positive curves use
$\sigma^{-1}=(a')^2\ell'$, where
$a'=a^{-1}$ and $\ell'=\ell^{-1}$. Divide by
\[
       (\ell,a)\longmapsto(q^2\ell,q^{-1}a),
                  \qquad 0<q<1 .                  \tag{C.EN.7}
\]
Every local root chart is the already verified
free finite quotient of an elliptic product in
(C.PE.8), so the resulting compact threefold
$\widetilde Y$ is smooth and proper over $V$.
Away from the four exceptional curves the
actual product coordinate is
$\ell=\sigma c^2,\ a=c^{-1}$, with
$[c]\in E=\C^*/q^{\Z}$.

Dualization gives an involution of
$\widetilde Y$: it takes the global meromorphic
coordinates $(\ell,a)$ to $(\ell^{-1},a^{-1})$
over $h$. Equivalently, a positive chart maps
to the target negative chart with unchanged
numerical coordinates $(\ell',a')=(\ell,a)$.
It conjugates (C.EN.7) to its inverse, so
descends to the quotient and still squares
to one. In the product coordinate it is
\[
                       (v,[c])\longmapsto(hv,[c^{-1}]).
                                                               \tag{C.EN.8}
\]
It covers $t\mapsto-t$.

We identify the component over the strict
transform of $R_1$. The restriction of $L$
there is trivial; in a compatible frame,
\[
 \sigma|_{R_1}=\phi(z)
   =\frac{(z-1)(z-2)}{(z+1)(z+2)},\qquad
                    \phi(-z)=\phi(z)^{-1}.          \tag{C.EN.9}
\]
Let $C$ be the smooth elliptic curve
\[
 C:\quad Y^2=(z^2-1)(z^2-4),\qquad
 y=\frac{Y}{(z+1)(z+2)},\quad y^2=\phi(z).
                                                               \tag{C.EN.10}
\]
Its four simple branch values are $1,2,-1,-2$.
The exact finite product-cover map from UN18 is
\[
 C\times E\longrightarrow B,\qquad
 (z,y,[u])\longmapsto[\ell=u^2,\ a=y/u].
                                                               \tag{C.EN.11}
\]
It is \'etale, including at the four root values,
by the two Jacobians in (C.PE.13). Its deck
involution is
\[
 \eta:(z,y,[u])\longmapsto(z,-y,[-u]).
                                                               \tag{C.EN.12}
\]
The translation $[u]\mapsto[-u]$ is a nonzero
two-torsion translation, so $\eta$ is free.
Thus $B=(C\times E)/\langle\eta\rangle$ is
bielliptic. The involution (C.EN.8) lifts to
\[
 k:(z,y,[u])\longmapsto(-z,y^{-1},[u^{-1}])
       \quad\text{or}\quad
       (z,Y,[u])\longmapsto(-z,Y,[u^{-1}]).
                                                               \tag{C.EN.13}
\]
These formulas commute with $\eta$.

On $C$, the involution $j:(z,Y)\mapsto(z,-Y)$
and $k_C:(z,Y)\mapsto(-z,Y)$ each act by $-1$
on $dz/Y$. The latter has four fixed points,
the two points over each of $0,\infty$.
Their product $jk_C$ has no fixed point:
over zero its equation would require $Y=0$,
and at infinity it exchanges the two points.
It is consequently a nonzero two-torsion
translation of the elliptic curve $C$.
It follows that $\eta$ and $\eta k$ are
free on $C\times E$, whereas $k$ has exactly
$4\cdot4=16$ fixed points. They occur in
eight free $\eta$-orbits. At each, the
derivative of $k$ is $(-1,-1)$.

The quotient $B/\langle h\rangle$ therefore
has exactly eight $A_1$ singularities and
no branch curve. Its minimal resolution
$M$ is Enriques. One exact verification
uses the intermediate quotient
$(C\times E)/\langle k\rangle$. Its
minimal resolution is K3 by UN23.
The involution $\eta$ extends over its
sixteen exceptional curves, permuting
them in pairs. It is free away from
them because both $\eta$ and $\eta k$
are free upstairs, and free on them
because they are paired. Its quotient
is precisely $M$. Equivalently, the
canonical form has character $-1$
under $\eta$ and $+1$ under $k$;
its square descends through the
explicit crepant $A_1$ resolutions,
while $p_g(M)=q(M)=0$. Thus $K_M$
has exact order two and $e(M)=12$.
This is an actual noncyclic product
quotient, consistent with the two
Enriques holonomy cases of
\cite[Theorem~2.1 and Example~2.3]{Yoshihara1995}.

It remains to realize this surface resolution
inside the threefold, not just on an abstract
component. Near $R_0$ the four root curves
are absent, so (C.EN.8) is a product.
Its fixed locus over zero consists of four
disjoint copies of $R_0$, indexed by $E[2]$.
Near either node of $R_0+R_1$, choose a
local coordinate $w$ at the corresponding
two-torsion point of $E$, so inversion is
$w\mapsto-w$. The ambient action is
\[
 (u,v,w)\longmapsto(u,-v,-w),\qquad
                         T=u^2v^2 .
                                                               \tag{C.EN.14}
\]
Its quotient has invariant coordinates
$A_0=v^2,B_0=vw,C_0=w^2$, with
$A_0C_0=B_0^2$ and $T=u^2A_0$.
The two smooth resolution charts are
\[
 \begin{array}{lll}
 A_0=r,&B_0=rs,&C_0=rs^2,\quad T=u^2r,\\
 A_0=r's'^2,&B_0=r's',&C_0=r',\quad T=u^2r's'^2 .
 \end{array}                                             \tag{C.EN.15}
\]
They glue by $s'=s^{-1}$ and $r'=rs^2$.
Globally they resolve the four fixed
curves. Restricting to $u=0$ is exactly
the minimal $A_1$ resolution of the
eight surface singularities just found.
In particular the strict central component
is the smooth Enriques surface $M$.

The component over $R_0$ becomes
$A=R_0\times(E/\{\pm1\})\simeq\PP^1\times\PP^1$.
Each exceptional surface of (C.EN.15)
is the projective normal-direction bundle
of a fixed curve. That normal bundle is
$\OO_{R_0}(-2)\oplus\OO_{R_0}$: the
product trivialization is valid on an
entire neighbourhood of $R_0$, and
$R_0^2=-2$. Hence these surfaces are
$E_i\simeq\mathbb F_2$.
The four root Hopf components are
exchanged in two pairs and give
$Q_1,Q_2$. Squaring the base parameter,
and the exact monomials in (C.EN.15),
give every coefficient in (C.EN.1).
Those same monomials show that its
reduced support is simple normal
crossing.

The full incidence calculation is
\[
 \begin{array}{c|c}
 \text{intersection}&\text{reduced space}\\ \hline
 M\cap A&\PP^1\sqcup\PP^1\\
 M\cap Q_j&\text{one elliptic curve}\quad(j=1,2)\\
 A\cap E_i&\PP^1\quad(i=1,\ldots,4)\\
 M\cap E_i&\PP^1\sqcup\PP^1\quad(i=1,\ldots,4)\\
 M\cap A\cap E_i&\text{two points}\quad(i=1,\ldots,4).
 \end{array}                                             \tag{C.EN.16}
\]
There are no other nonempty intersections.
Thus inclusion--exclusion gives, with
every nonzero term displayed,
\[
 e(S)=12+4+4\cdot4-(4+4\cdot2+4\cdot4)+8=12 .
                                                               \tag{C.EN.17}
\]

For the compact construction, the base involution
also fixes the smooth elliptic fibre at infinity.
The ambient fixed loci there are four elliptic
curves, and the same transverse $A_1$ charts
resolve them. Their exceptional surfaces and
all intersections are elliptic ruled and
have Euler number zero. At the remaining
nodal base values the elliptic-product
coordinate gives a nodal curve times $E$,
also of Euler number zero. The quotient
and the resolutions are proper and compact.
Over a general $T$, the two $t$-sheets
are exchanged by $h$, so their quotient
is exactly one product of two elliptic
curves. The Stein factor has generic
degree one over the normal curve
$\PP^1$, proving connectedness.

Finally perform six of the point-then-quartic
operations of UN4, all over $T=0$ and
away from $M$. They can begin at fresh
smooth points of an $E_i$ outside its
intersections. Their exceptional plane
quartics and subsequent centres remain
disjoint from $M$. Each changes the
reduced Euler number by $-2$, preserves
the smooth total space and the normal
crossing support, and leaves the
punctured family unchanged. The resulting
support has Euler number zero, with
the same smooth Enriques component.
Its nonzero even pluricanonical form
precludes a covering family of rational
curves, so it is non-uniruled.
The exact inverse-map resolution argument
of UN14 proves essentiality. This completes
the compact occurrence construction.
\input{sections/enriques_support.inc.tex}
\end{proof}
